% demo-mypseudocode.tex -- mypseudocode / mypseudocode*
\documentclass[sectioning=nobreak]{myreport}
\usepackage{clrscode3e}   % provides the codebox environment (required)
\usepackage{mydemo}
\begin{document}
\demoheader{mypseudocode}{Pseudocode in the style of Cormen et al.
(clrscode3e), numbered together with the programs.}

\mysection*{Syntax}
\begin{demosyntax}
\usepackage{clrscode3e}   % in the preamble: required
\begin{mypseudocode}{<label>}[<short caption>]{<caption>}[<position>]
  ... clrscode3e commands ...
\end{mypseudocode}
\begin{mypseudocode*}{<label>}{<caption>}[<position>]
\end{demosyntax}
\demonote{The label is \texttt{code:pseudo:<label>}. The environment uses
\texttt{codebox} from \texttt{clrscode3e}, which you must load yourself;
it cannot be combined with \texttt{algpseudocode} (both define
\texttt{\textbackslash For}, \texttt{\textbackslash If}, \dots).}

\mysection*{Example}
\begin{demo}
\begin{mypseudocode}{isort}{Insertion sort.}[H]
\proc{Insertion-Sort}(A)
\li \For $j \gets 2$ \To $\attrib{A}{length}$
\li   \Do $\id{key} \gets A[j]$
\li       $i \gets j-1$
\li       \While $i > 0$ and $A[i] > \id{key}$
\li         \Do $A[i+1] \gets A[i]$
\li             $i \gets i-1$
            \End
\li       $A[i+1] \gets \id{key}$
      \End
\end{mypseudocode}
See \myref{code:pseudo:isort}.
\end{demo}
\end{document}
